% ══════════════════════════════════════════════════════════════════════
% § 6  Discussion
% ══════════════════════════════════════════════════════════════════════

\subsection{Summary}

We have classified all four-dimensional $\N=1$ asymptotically free
$\SU(N)\times\SU(N)$ quiver gauge theories in the equal-rank,
non-Veneziano large-$N$ sector, using symbolic $a$-maximisation to
compute IR R-charges and central charges.  Nine universality classes
emerge; every class with a defined leading-$N$ $a$-maximum satisfies
$a=c$ to leading order in $1/N$.

Applying the boundary analysis of \cite{IW2005}, we have computed the
boundary one-loop coefficients $b_{0}^{(a)}(g_{b}^{\star})$ for every
class and classified each by the sign pattern
$(\text{sign}\,\alpha_{1},\text{sign}\,\alpha_{2})$ of the
leading-$N$ coefficients.  Class 2 realises the Figure~5 topology;
classes $4,16,21,23$ realise a mixed topology in which one axis SCFT
attracts generic flows; class 44 is leading-$N$ marginal at both
boundaries.  Classes $1$, $9$, and $877$ together with the 20
flat-direction theories excluded from the nine-class tabulation form
a single \emph{nonSCFT} umbrella: in each case the leading-$N$
$a$-maximisation does not identify a unitary interacting SCFT, either
because it returns R-charges below the unitarity bound (classes 1, 9,
877 — at $R=0$, $R=1/2$, and $R=0$ on surviving matter respectively)
or because the critical point is not a genuine local maximum (flat
direction).  Class 877 is further singled out by its empty-node
morphology: the 9 surviving representatives admit a unique interior
$a$-maximum at $a/N^{2}=0$, but the sub-quiver decoupling limit
required for the boundary analysis is subtle for an all-bifundamental
matter sector and is left for a separate treatment.

Our main result is that the strictly non-Veneziano landscape contains
no Figure~6 realisation, but that a simple single-parameter Veneziano
deformation of class~16 (adding $\Nf^{(2)}$ fundamental pairs to
node~2 in the window $0.427\,N-1<\Nf^{(2)}<N-1$) produces precisely
the both-axes-stable RG topology of \cite[Fig.~6]{IW2005}.  The
interior $a$-maximal critical point persists as a saddle inside this
window, and the two axis SCFTs become the IR attractors with basins
separated by the stable manifold of the saddle.  This is, to our
knowledge, the first explicit realisation of this RG topology.

\subsection{Implications}

\paragraph{$a=c$ in the equal-rank non-Veneziano sector.}  The
uniform appearance of $a=c$ across all classes with a defined
leading-$N$ central charge suggests the non-Veneziano restriction is
itself the relevant filter: the
subleading, flavour-asymmetric contribution to $\Tr R^{3}-\Tr R$ that
would distinguish $a$ from $c$ is generically the
$\mathcal O(N)$-many fundamentals' contribution, and setting
$\Nf=\mathcal O(N^{0})$ removes it at leading order.  Quantifying the
sub-leading $a-c$ split as a function of a Veneziano parameter
$x^{(a)}=\Nf^{(a)}/N$ is a natural extension.

\paragraph{Holographic dual of the Figure~6 phase.}  In the Figure~5
topology, the single interior SCFT that attracts generic flows
admits a holographic interpretation through a single interior AdS
throat.  The Figure~6 topology, by contrast, has two distinct IR
attractors on the axes; a holographic dual would require a bulk
geometry with two separate IR asymptotic regions, connected in the
UV, selected by the UV Wilsonian initial condition.  Whether the
deformed class~16 realises such a bulk dual --- and whether the
separatrix structure admits a gravity interpretation --- is an open
question.

\paragraph{Higher-loop robustness.}  The boundary criterion used here
is one-loop via NSVZ and linear in the bifundamental anomalous
dimension.  Two- and higher-loop corrections can shift the
sign-flip thresholds by $\mathcal{O}(1/N)$ amounts.  The lower edge
of the class-16 Figure~6 window at $\Nf^{(2)}=0.427\,N-1$ is
therefore expected to receive $\mathcal{O}(1/N)$ corrections; the
upper edge $\Nf^{(2)}=N-1$ is the UV-AF cliff and is exact.  Because
the boundary coefficient $\alpha_{2}=-0.5730$ on axis 2 is
$\mathcal{O}(N)$ and of definite sign throughout the window, the
qualitative Figure~6 phase assignment is robust against higher-loop
corrections.

\paragraph{Boundary-versus-interior $R_{\text{bif}}$.}  We have used
the interior $R_{\text{bif}}$ (from full-theory $a$-maximisation) as
a leading-$N$ proxy for the boundary sub-theory $R_{\text{bif}}$.
The two prescriptions agree qualitatively (same sign pattern) under
the mild assumption that neither has $R_{\text{bif}}$ cross the
free-field value $2/3$, which holds for all classes with a defined
interior $a$-maximum.  A
refinement using the strict sub-theory $a$-max for the boundary
R-charges would change the quantitative value $0.427$ at the
Figure~6 lower edge by $\mathcal O(1/N)$ but not alter the existence
or the width-scaling $\approx 0.573\,N$ of the window.

\subsection{Future work}
\label{sec:future}

\begin{itemize}
  \item \textbf{$n$-node quivers.}  Extending this analysis to linear
    and cyclic quivers with $n\ge 3$ nodes requires a
    $2^{n}$-dimensional sign-pattern analysis on the boundaries; we
    expect richer topologies including partial basins and sub-quiver
    attractors.
  \item \textbf{$\SO/\Sp$ sectors.}  The database also contains
    $\SU$--$\SO$, $\SU$--$\Sp$ and $\SO$--$\Sp$ two-node theories.  A
    parallel analysis of these mixed-group sectors in the
    non-Veneziano limit should uncover additional realisations of the
    three topologies and possibly new Figure~6 windows.
  \item \textbf{Sub-leading treatment of class~44 and of the nonSCFT
    umbrella.}
    Class 44's boundary coefficients vanish at leading $N$; the
    subleading structure would determine its RG topology exactly.
    The nonSCFT umbrella — classes 1, 9, 877 and the 20
    flat-direction theories — is where the leading-$N$ $a$-maximisation
    fails to identify a unitary interacting SCFT: classes~1 and~9
    have R-charges below unitarity ($R=0$ and $R=1/2$), class~877 has
    $a/N^{2}=0$ with all surviving matter R-charges vanishing and
    additionally requires a separate sub-quiver decoupling treatment
    because its empty-node morphology places all matter on
    bifundamentals, and the flat-direction theories have
    non-uniquely-determined IR R-charges.  Lifting any of these into
    the interior of the polytope by sub-leading corrections (or
    finding a different R-symmetry extremum) would give their
    physical IR.
  \item \textbf{Other mixed-phase Figure~6 candidates.}  Class 21
    admits a narrow window $0.6754\,N+2<\Nf^{(1)}<N+2$ that may
    give a second Figure~6 realisation under analogous deformation;
    because class 21's node 2 is leading-$N$ marginal, a careful
    sub-leading analysis is needed to confirm this.
  \item \textbf{Explicit two-coupling RG numerics.}  Numerical
    integration of the two-coupling NSVZ flow in the
    $(g_{1},g_{2})$-plane, with R-charges updated self-consistently
    along the trajectory, would make the separatrix and basin
    structure in the class-16 Figure~6 window fully manifest.
  \item \textbf{Veneziano phase diagram.}  The sign-flip
    $\alpha_{1}<0\to\alpha_{1}>0$ at $\Nf^{(2)}=0.427\,N-1$ is the
    boundary of the Figure~6 phase in the $(\Nf^{(1)},\Nf^{(2)})$
    plane for class~16.  Mapping out the full two-dimensional phase
    diagram in $(x^{(1)},x^{(2)})$ Veneziano parameters would
    generalise the present results to the Veneziano regime and
    naturally connect to the full database.
\end{itemize}

\paragraph{Acknowledgements.}  We thank the organisers of [...] for
useful discussions.  Computations were performed with the symbolic
$a$-maximisation pipeline described in the companion code repository.
